\documentclass[10pt,a4paper]{article}
\usepackage{arabtex}
\usepackage[OT1,T1,LFE,LAE]{fontenc}
\usepackage[utf8]{inputenc}
\usepackage[arabic,english,farsi]{babel}
\usepackage{amsmath,amsfonts} % Math packages
\usepackage{amssymb}
%\usepackage{cmap}

\usepackage{multicol}

\usepackage{graphicx}
\usepackage[caption=false]{subfig}
\usepackage{color}
\usepackage{float}
\usepackage{sidecap}
%\sidecaptionvpos{figure}{c}
\usepackage{anysize}
\marginsize{2cm}{2cm}{2cm}{2cm}

\usepackage{listings}

\usepackage{appendix}
%\renewcommand{\appendixname}{Apéndices}
%\renewcommand{\appendixtocname}{Apéndices}
%\renewcommand{\appendixpagename}{Apéndices} 

\usepackage[colorlinks=true,plainpages=true,citecolor=blue,linkcolor=blue,urlcolor=cyan]{hyperref}
%\usepackage{hyperref}


%%% Equation and float numbering
\numberwithin{equation}{section}
\numberwithin{figure}{section}
\numberwithin{table}{section}


\newcommand{\horrule}[1]{\rule{\linewidth}{#1}}    % Horizontal rule

\newcommand{\titleText}{Network management and security \\ Laboratory Manual}

\title{
\normalsize In the name of Allah\\
\vspace{10pt}
\LARGE\FR{بسم \allah الرحمن الرحیم}
\vspace{10pt}
\begin{center}
    %	\newcommand{\HRule}{\rule{\linewidth}{0.5mm}}
    \begin{minipage}{0.48\textwidth}
        \begin{flushleft}
            \includegraphics[height=64pt,width=64pt]{../img/logo.png}
        \end{flushleft}
    \end{minipage}
    \begin{minipage}{0.48\textwidth}
        \begin{flushright}
            \includegraphics[height=64pt]{../img/eng-logo.png}
        \end{flushright}
    \end{minipage}
\end{center}
\vspace*{-64pt}
%	\horrule{0.5pt} \\[0.4cm]
\huge \titleText\\
\vspace{40pt}
%	\horrule{2pt} \\[0.5cm]
}
\author{
\huge University of Tehran\\
\LARGE \FR{دانشگاه تهران}\\
\\
\LARGE School of Electrical and Computer Engineering\\
\FR{دانشکده مهندسی برق و کامپیوتر}\\
\\
\Large Computer Network Lab\\
\FR{آزمایشگاه شبکه‌های کامپیوتری}\\
\\
\\
\\
\normalfont
Dr. Ahmad Khonsari - \FR{احمد خونساری}\\
\href{mailto:a_khonsari@ut.ac.ir}{a\_khonsari@ut.ac.ir}\\
\\
\normalsize
Amir Haji Ali Khamseh'i - \FR{امیر حاجی علی خمسه‌ء}\\
\href{mailto:khamse@ut.ac.ir}{khamse@ut.ac.ir}\\
\\
\normalsize
Sina Kashi pazha - \FR{سینا کاشی پزها}\\
\href{mailto:sina\_kashipazha@ut.ac.ir}{sina\_kashipazha@ut.ac.ir}\\
\\
\normalsize
Mohammad Ali Shahsavand - \FR{محمد علی شاهسوند}\\
\href{mailto:mashahsavand@ut.ac.ir}{mashahsavand@ut.ac.ir}\\
\\
\normalsize
Amirahmad Khordadi - \FR{امیر احمد خردادی}\\
\href{mailto:a.a.khordadi@ut.ac.ir}{a.a.khordadi@ut.ac.ir}
}

\date{\vspace{30pt}\today\\\vspace{10pt}{\selectlanguage{farsi}\today}}

\usepackage{fancyhdr}
\pagestyle{fancy}
%\pagestyle{fancyplain}
\fancyhf{}
\fancyhead[L]{\footnotesize Computer Network Lab \\ \FR{آزمایشگاه شبکه‌های کامپیوتری}}
\fancyhead[R]{\footnotesize \titleText}
\fancyfoot[R]{\footnotesize School of Electrical and Computer Engineering\\\FR{دانشکده مهندسی برق و کامپیوتر}}
\fancyfoot[C]{\thepage}
\fancyfoot[L]{\footnotesize University of Tehran \\ \FR{دانشگاه تهران}}
\renewcommand{\footrulewidth}{0.8pt}
\renewcommand{\headrulewidth}{1pt}            % Remove header underlines
\renewcommand{\footrulewidth}{1pt}                % Remove footer underlines
\setlength{\headheight}{13.6pt}


\begin{document}
    \selectlanguage{english}

    %\vspace*{-1.5cm}
    \maketitle

    %\tableofcontents

    \pagebreak

\section{Exercises on secure applications}
\subsection{Excercise One}
In this excercise we will study security vulnerability of ftp and telnet protocol . To do so, wi will create mininet topology with single hub\footnote{hub forwards incoming packets to all of its ports, which means it always floods packets} connected to three hosts then we will connect from h1 to h2 through ftp and telnet connection and capture h1 password on h3. Let's do it.

\begin{itemize}
	\setlength{\itemindent}{10pt}
	\item Start pox controller with below command to force mininet switches act like hub.
	\begin{enumerate}
		\setlength{\itemindent}{10pt}
		\item [\$]  python pox.py opeenflow.of\_01 {-}{-}address=127.0.0.1 {-}{-}port=6337 forwarding.hub
	\end{enumerate}
	\item Run below command to start mininet with one single switch and three hosts and connect it to pox controller.
	\begin{enumerate}
		\setlength{\itemindent}{10pt}
		\item [\$]  sudo mn {-}{-}topo single,3 {-}{-}controller remote,ip=127.0.0.1,port=6633
	\end{enumerate}
	\item start ftp server on h2 with :
	\begin{enumerate}
		\setlength{\itemindent}{10pt}
		\item [\#] /usr/sbin/vsftpd  
	\end{enumerate}
	\item Run \textbf{wireshark \&} on h3
	\item Login to h2 ftp from h1
	\item Capture h1 password on wireshark output
	
\end{itemize}

Repeat the above experiment, but use telnet to connect from h1 to h2 and capture h1 password on h3. \footnote{Don't foget to restart xinetd with `/etc/init.d/xinetd restart` on h2 to start telnet server. }

\subsubsection*{ Lab Report}
\begin{itemize}
	\setlength{\itemindent}{0pt}
	\item Can you see the login ID and the password in the FTP experiment? Submit the two packets you captured.
	\item Can you see the login ID and the password in the TELNET experiment? Submit the packets you captured.
	\item What is the difference between FTP and TELNET in their transmission of user ID’s and passwords? Which one is more secure?
\end{itemize}

\subsection{Excercise Two}


\subsubsection*{ Lab Report}
\begin{itemize}
	\setlength{\itemindent}{0pt}
	\item In each experiment, can you extract the password from the tcpdump output? Can you read the IP, TCP, SSH headers? Can you read the TCP data?
	\item What is the client protocol (and version) used in both cases?
	\item What is the port number used by the ssh server? What is the port number used by the sftp server? Justify your answer using the wireshark output and the /etc/services file.
\end{itemize}

    \pagebreak

\section{Exercises on Firewalls and Iptables}
\subsection*{Excercise Three}
Start mininet with default topology and Execute iptables -L -v on h1 and h2 to list the existing rules in the filter table. Save the output for the lab report.\\
\\
 Append a rule to the end of the INPUT chain, by executing

\begin{itemize}
	\setlength{\itemindent}{10pt}
	\item [h2>] \textbf{iptables -A INPUT -v -p TCP {-}{-}dport 23 -j DROP} 
\end{itemize}

\setlength{\parindent}{0pt}
on h2. Run iptables -L -v again on both hosts to display the filter table. Save the output.\\


\begin{itemize}
	\setlength{\itemindent}{0pt}
	\item Start telnet server on h2 with `/etc/init.d/xinetd restart`.
	\item Capture packets on both hosts with wireshark
	\item Try to login with telnet from h1 to h2
\end{itemize}

\subsubsection*{ Lab Report}
\begin{itemize}
	\setlength{\itemindent}{0pt}
	\item Can you telnet to the host from the remote machine?
	\item From the wireshark output, how many retries did telnet make? Explain the exponential backoff algorithm of TCP timeout and retransmission.
\end{itemize}

\subsection*{Excercise Four}
Keep previous mininet running and delete the rule created in the last exercise on h2, by:

\begin{itemize}
	\setlength{\itemindent}{10pt}
	\item [h2>] \textbf{iptables -D INPUT -v -p TCP {-}{-}dport 23 -j DROP} 
\end{itemize}

\setlength{\parindent}{0pt}
Then, append a new rule to the INPUT chain:

\begin{itemize}
	\setlength{\itemindent}{10pt}
	\item [h2>] \textbf{iptables -A INPUT -v -p TCP {-}{-}dport 23 -j REJECT {-}{-}reject-with tcp-reset} 
\end{itemize}

\setlength{\parindent}{0pt}
Execute \textbf{iptables -L -v} to display the new rule. On both machines in your topology, restart wireshark output, and then telnet from h1 to h2. Save the wireshark output for the lab report.

\subsubsection*{ Lab Report}
\begin{itemize}
	\setlength{\itemindent}{0pt}
	\item Explain the difference between the wireshark outputs of this exercise and the previous exercise. How many attempts did TCP make this time?
\end{itemize}


\end{document}
